\begin{lemma}[The sum rule in the common-time direction]
Let $F$ be a finite-dimensional real inner-product space, let $N$ be a natural
number, let $\Omega\subseteq\ell^2(F\times\R)$ and let $v$ be a viscosity
subsolution of $F[v]=0$ on $\Omega$ in the sense of \noderef{IsLimitSubOn} (so
in particular $\Omega$ is open and $v$ is continuous on $\Omega$). Let
$\theta_i>0$ for $1\le i\le N$, write
$P=\ell^2\bigl(\bigoplus_{i=1}^N\ell^2(F\times\R)\bigr)$ for the $\ell^2$-product
of $N$ copies of space-time, with elements $W=(w_i)_{i=1}^N$, and let
$\Psi:P\to\R$. Let $Z=(z_i)_{i=1}^N\in P$ with $z_i\in\Omega$ for every $i$,
assume $\Psi$ is $C^2$ near $Z$, and assume that
\[
  W\longmapsto\sum_{i=1}^N\theta_i\,v(w_i)-\Psi(W)
\]
has a local maximum at $Z$. Put
\[
  q_i=\theta_i^{-1}\bigl(\nabla\Psi(Z)_i\bigr)_x\in F,
  \qquad
  \zeta=\bigl((q_i,1)\bigr)_{i=1}^N\in P,
\]
where $\nabla\Psi(Z)_i\in\ell^2(F\times\R)$ is the $i$-th block of
$\nabla\Psi(Z)$ and $(\cdot)_x$ is its spatial part. Then
\[
  D^2\Psi(Z)[\zeta]\ge0 .
\]
\end{lemma}
\begin{proof}
Write $A=D^2\Psi(Z)$, regarded as a continuous bilinear form on $P$, and note
that neither $A$ nor $\zeta$ depends on the parameter $\eta$ introduced below:
both are determined by $\Psi$ and $Z$ alone.

\emph{Step 1: the sum rule at a fixed $\eta$.}
Fix $\eta>0$. The data
\[
  N,\qquad O_i=\Omega,\qquad g_i=\theta_i v,\qquad \Psi,\qquad Z,\qquad \eta
\]
satisfy every hypothesis of \noderef{SumRuleJets}: each $O_i=\Omega$ is open and
each $g_i=\theta_i v$ is continuous on it, both by \noderef{IsLimitSubOn}
(continuity of $\theta_i v$ being continuity of $v$ multiplied by the constant
$\theta_i$); $z_i\in\Omega$ for every $i$ by hypothesis; $\Psi$ is $C^2$ near
$Z$; the function $W\mapsto\sum_i g_i(w_i)-\Psi(W)$ has a local maximum at $Z$
by hypothesis, since $g_i(w_i)=\theta_iv(w_i)$; and $\eta>0$. That lemma
therefore supplies symmetric continuous bilinear forms $X_i$ on
$\ell^2(F\times\R)$ with
\begin{equation}\label{eq:srctd-jet}
  \bigl(\nabla\Psi(Z)_i,\ X_i\bigr)
   \in\overline J^{2,+}\bigl(\theta_i v\bigr)(z_i)
  \qquad(1\le i\le N)
\end{equation}
and
\begin{equation}\label{eq:srctd-matrix}
  \sum_{i=1}^NX_i[\xi_i]\le A[\xi]+\eta\,\bigl|A\xi\bigr|^2
  \qquad\bigl(\xi=(\xi_i)_i\in P\bigr),
\end{equation}
where $|A\xi|$ is the norm of the linear functional $A(\xi,\cdot)$.

\emph{Step 2: dividing each jet by $\theta_i$.}
Fix $i$. Since $\theta_i>0$, \noderef{ClosedSuperjetScale} applied with
$\theta=\theta_i$, $g=v$, $z=z_i$, $p=\nabla\Psi(Z)_i$ and $X=X_i$ turns
\eqref{eq:srctd-jet} into
\begin{equation}\label{eq:srctd-vjet}
  \bigl(\theta_i^{-1}\nabla\Psi(Z)_i,\ \theta_i^{-1}X_i\bigr)
   \in\overline J^{2,+}v(z_i).
\end{equation}
The spatial part of $\theta_i^{-1}\nabla\Psi(Z)_i$ is
$\theta_i^{-1}\bigl(\nabla\Psi(Z)_i\bigr)_x=q_i$, scalar multiplication on
$\ell^2(F\times\R)$ being componentwise.

\emph{Step 3: the viscosity inequality on each rescaled jet.}
Since $z_i\in\Omega$ and \eqref{eq:srctd-vjet} is a closed superjet of $v$ at
$z_i$, \noderef{LimitSubClosedSuperjet} applies to it and gives
\[
  0\le\bigl(\theta_i^{-1}X_i\bigr)\bigl[(q_i,1)\bigr],
\]
the direction being formed from the spatial part $q_i$ of the jet's gradient and
the time component $1$, exactly as that lemma prescribes. Now
$\bigl(\theta_i^{-1}X_i\bigr)[(q_i,1)]=\theta_i^{-1}X_i[(q_i,1)]$, so
multiplying by $\theta_i>0$ gives
\begin{equation}\label{eq:srctd-each}
  X_i\bigl[(q_i,1)\bigr]\ge0\qquad(1\le i\le N).
\end{equation}

\emph{Step 4: evaluating the matrix inequality on $\zeta$.}
The $i$-th block of $\zeta$ is $(q_i,1)$, so taking $\xi=\zeta$ in
\eqref{eq:srctd-matrix} and summing \eqref{eq:srctd-each} over $i$ gives
\begin{equation}\label{eq:srctd-eta}
  0\le\sum_{i=1}^NX_i\bigl[(q_i,1)\bigr]\le A[\zeta]+\eta\,|A\zeta|^2 .
\end{equation}

\emph{Step 5: letting $\eta\downarrow0$.}
Step 1 was carried out for an arbitrary $\eta>0$, and, as recorded at the start
of the proof, neither $A$ nor $\zeta$ depends on $\eta$. So \eqref{eq:srctd-eta}
says
\[
  0\le A[\zeta]+\eta\,|A\zeta|^2\qquad\text{for every }\eta>0 .
\]
Suppose, for contradiction, that $A[\zeta]<0$, and put
\[
  \eta_\star=\frac{-A[\zeta]}{2\bigl(|A\zeta|^2+1\bigr)}>0,
\]
which is positive because the numerator is positive and the denominator is at
least $2$. Then
\[
  \eta_\star|A\zeta|^2\le\eta_\star\bigl(|A\zeta|^2+1\bigr)=\frac{-A[\zeta]}2,
\]
so
\[
  A[\zeta]+\eta_\star|A\zeta|^2\le A[\zeta]-\frac{A[\zeta]}2
   =\frac{A[\zeta]}2<0,
\]
contradicting the displayed inequality at $\eta=\eta_\star$. Hence
$A[\zeta]\ge0$, which is the assertion.
\end{proof}
